\documentclass[12pt]{article}
\usepackage[utf8]{inputenc}
\usepackage[T1]{fontenc}
\usepackage{hyperref}
\title{Takua Language Resources for K-12 Curriculum Development}
\author{Compiled by Claude Sonnet 5 (claude-sonnet-5) for Andr\'as Kornai}
\date{}
\begin{document}
\maketitle

\section{Introduction}
Takua (ISO 639-3 tkz) is a North Bahnaric Austroasiatic language spoken in the mountains of Qu\h{a}ng Nam and Qu\h{a}ng Ng\~ai provinces, Vietnam. Estimates are about 6,000 speakers (Wikipedia, 2007) and 6,800 to 6,900 people (Joshua Project). It should not be confused with Katua, a different language. Only encyclopedic and missions-database pages appeared in the search results.

\section{Primary Educational Resources}
None found.

\section{Linguistic and Reference Materials}
\begin{itemize}
\item Takua language (Wikipedia): \url{https://en.wikipedia.org/wiki/Takua_language}
\item Takua people (Wikipedia): \url{https://en.wikipedia.org/wiki/Takua_people}
\end{itemize}
No dictionary, grammar or orthography was found.

\section{Religious and Community Resources}
Joshua Project profile of the Takua people (about a quarter Christian, per that source): \url{https://joshuaproject.net/people_groups/15208}\\
Joshua Project language page: \url{https://joshuaproject.net/languages/tkz}

\section{Digital and Audio Resources}
Nothing found.

\section{Curriculum Development Context}
Little to nothing exists for classroom use. Curriculum work would require new documentation, ideally building on materials for related Bahnaric languages and with community input.
\end{document}